\section{Stability}
\label{sec:stability}

Stability is a property of the equilibrium points of a dynamical system, not of the system as a
whole. Adding a feedback law turns the controlled system $x_{k+1} = f(x_k, u_k)$ into the
autonomous system $x_{k+1} = f(x_k, \kappa(x_k)) = \tilde f(x_k)$, and two questions arise: can
an unstable equilibrium be made stable by the controller, and can an equilibrium that was
stable be made unstable by it? For MPC the second question is the one that needs an answer,
because $\kappa_N$ is a nonlinear law that no classical criterion covers.

\subsection{Definitions}

Consider an autonomous system $x_{k+1} = \tilde f(x_k)$ with an equilibrium at the origin,
$\tilde f(0) = 0$.

\begin{definition}[stability]
The origin is \emph{stable} if for every $\varepsilon > 0$ there exists $\delta(\varepsilon) > 0$
such that $\norm{x_{k_0}} \le \delta$ implies $\norm{x_k} \le \varepsilon$ for all $k \ge k_0$.
It is \emph{asymptotically stable} if it is stable and, in addition, $x_k \to 0$ for every
initial state in some neighbourhood of the origin, the \emph{region of attraction}. It is
\emph{globally asymptotically stable} if the region of attraction is the whole state space.
\end{definition}

Stability says that small initial perturbations produce trajectories that remain forever close
(\cref{fig:epsdelta}); asymptotic stability adds that they return. Both are local notions
unless stated otherwise.

\begin{figure}[t]
\centering
\begin{tikzpicture}[x=1cm,y=1cm,>=Stealth]
  \draw[thick,sap,dashed] (0,0) circle (2.0);
  \draw[thick,mpc] (0,0) circle (0.8);
  \fill (0,0) circle (0.05);
  \node[below right,font=\small] at (0,0) {$x^{eq}$};
  \node[mpc,font=\small] at (0.95,-0.95) {$\delta$};
  \node[sap,font=\small] at (1.6,1.6) {$\varepsilon$};
  \fill[mpc] (-0.5,0.4) circle (0.05);
  \node[mpc,font=\small,left] at (-0.5,0.4) {$x_{k_0}$};
  \draw[thick,->] plot[smooth] coordinates {(-0.5,0.4) (0.3,1.3) (1.4,0.9) (1.1,-0.6) (-0.4,-1.2) (-1.3,-0.2) (-0.6,0.9) (0.5,0.6) (0.6,-0.3)};
\end{tikzpicture}
\caption{The $\varepsilon$--$\delta$ definition of stability: a trajectory that starts within
$\delta$ of the equilibrium never leaves the ball of radius $\varepsilon$.}
\label{fig:epsdelta}
\end{figure}

\subsection{Linear systems}

For $x_{k+1} = A x_k$ the answer is complete. The solution is $x_k = A^k x_0$, and everything
depends on the eigenvalues of $A$: the origin is asymptotically stable if and only if all the
eigenvalues lie strictly inside the unit circle; it is unstable if at least one lies outside;
eigenvalues on the circle give simple stability when they are simple, and need a separate check
when they are repeated. This criterion is not enough for MPC: the model may be nonlinear, and
even for a linear model the controller $\kappa_N$ is nonlinear.

\subsection{Lyapunov functions}

\begin{definition}[definite functions]
A scalar function of the state $V: \R^n \to \R$ is \emph{positive definite} if $V(0) = 0$ and
$V(x) > 0$ for every $x \ne 0$; it is \emph{positive semi-definite} if $V(0) = 0$ and $V(x) \ge
0$ for every $x$. It is negative definite, or semi-definite, when $-V$ is positive definite, or
semi-definite. The \emph{local} versions of these definitions require the inequalities only in a
neighbourhood of the equilibrium.
\end{definition}

The quadratic forms $V(x) = x\T P x$ are the standard example: such a form is positive definite
if and only if all the eigenvalues of $P$ are positive. Not every candidate works: $V(x) =
2\sin^2(x_1) + 3x_2^2$ is positive definite only locally, since it vanishes at $(\pi, 0)$. The level sets
$\Omega_c = \{x : V(x) = c\}$ of a positive definite function are closed curves around the
origin, nested as $c$ grows; they are the geometric picture behind the theorems that follow.

The \emph{increment} of $V$ along the trajectories of the system is
\begin{equation}
  \Delta V(x) = V\big(\tilde f(x)\big) - V(x),
\end{equation}
the change of $V$ from one state to the next. It is computed along the trajectories, not in an
arbitrary direction.

\begin{theorem}[Lyapunov]
Let the origin be an equilibrium of $x_{k+1} = \tilde f(x_k)$. If there exists a continuous
function $V$, positive definite at least locally around the origin, such that $\Delta V$ is negative
semi-definite at least locally, then the origin is stable. If, in addition, $\Delta V$ is
negative definite, then the origin is asymptotically stable.
\end{theorem}

The condition $\Delta V \le 0$ is often written as $\Delta V(x) \le -h(x)$ with $h \ge 0$,
which says \emph{how much} $V$ must decrease at each step; this form is the one used in the
next section. The intuition is in \cref{fig:lyapunov}: if $V$ is positive and decreases along
the trajectory, the trajectory is confined to level sets of smaller and smaller value, and if
the decrease is strict then $V \to 0$ and with it $x_k \to 0$. A function tried as a candidate is
a \emph{candidate Lyapunov function}; one that satisfies both conditions is a \emph{Lyapunov
function}. The theorem is a sufficient condition: failing to find a Lyapunov function proves
nothing.

\begin{figure}[t]
\centering
\begin{tikzpicture}[x=1cm,y=1cm,>=Stealth]
  \foreach \r in {0.5,1.0,1.5,2.0} {
    \draw[mpc!60] plot[smooth cycle] coordinates {(\r,0) (0.7*\r,0.8*\r) (-0.2*\r,1.1*\r) (-1.1*\r,0.3*\r) (-0.8*\r,-0.8*\r) (0.3*\r,-1.0*\r)};
  }
  \node[mpc,font=\small] at (2.45,0.0) {$\Omega_{c}$};
  \fill (0,0) circle (0.05);
  \draw[thick,->] plot[smooth] coordinates {(1.9,0.9) (0.9,1.4) (-0.9,0.9) (-1.1,-0.3) (-0.2,-0.9) (0.8,-0.4) (0.5,0.4) (-0.2,0.3) (-0.1,-0.1)};
  \fill[sap] (1.9,0.9) circle (0.06);
  \node[sap,font=\small,right] at (1.9,0.9) {$x_0$};
\end{tikzpicture}
\caption{Level sets of a positive definite function and a trajectory along which $V$ decreases
at every step: each level set traps the trajectory inside a smaller one.}
\label{fig:lyapunov}
\end{figure}

\subsection{Stability of MPC}

Take the regulation problem: the system has an equilibrium at the origin, $f(0,0) = 0$, the
stage cost $\ell(x,u)$ is positive definite in $(x,u)$, and the terminal cost $V_f(x)$ is
positive definite in $x$. These assumptions are satisfied by any reasonable regulation cost,
such as \eqref{eq:quadcost} with $Q$ and $R$ positive definite.

\paragraph{The equilibrium survives.} The closed loop is $x_{k+1} = f(x_k, \kappa_N(x_k)) =
\tilde f(x_k)$, so the origin is an equilibrium of $\tilde f$ if and only if $\kappa_N(0) = 0$. It
is: from $x_k = 0$ the sequence of all zeros is admissible, keeps the state at zero and costs
nothing, and since the stage and terminal costs are positive definite no other sequence can
cost less. The optimizer returns zero, and the controller does not move the equilibrium.

\paragraph{The candidate.} The candidate Lyapunov function is the \emph{optimal cost of the MPC
problem}, $V^*_N(x_k) = V_N(x_k, \bu^*_k)$: a function of the state alone, computed by the
controller at every step in any case. It is positive definite: $V^*_N(0) = 0$ by the argument
above, and for $x \ne 0$ the first stage cost $\ell(x_k, u^*_{\ik{0}})$ alone is strictly
positive.

\paragraph{The decrease.} We need $V^*_N(x_{k+1}) - V^*_N(x_k) \le 0$, and $V^*_N(x_{k+1})$ is
the result of an optimization that we have not performed. The way around it is a bound: the
optimal cost at $k+1$ is never larger than the cost of \emph{any} admissible sequence at $k+1$,
so it is enough to exhibit one admissible sequence whose cost is already below $V^*_N(x_k)$. The
candidate is the one of \cref{sec:feasibility}: the tail of the optimal sequence of time $k$,
completed with the terminal control, which is admissible under the nominal assumption
$x_{k+1} = x^*_{\ik{1}}$ and condition \eqref{eq:terminalinv}. Writing both costs term by term,
the stage costs from $1$ to $N-1$ cancel and what is left is
\begin{equation}
\begin{split}
  V_N(x_{k+1}, \tilde\bu_{k+1}) - V^*_N(x_k)
  = {}& \underbrace{-\,\ell\big(x_k, u^*_{\ik{0}}\big)}_{\le 0} \\
  &+ \Big[\ell\big(x^*_{\ik{N}}, \kappa_f(x^*_{\ik{N}})\big)
       + V_f\big(f(x^*_{\ik{N}}, \kappa_f(x^*_{\ik{N}}))\big) - V_f\big(x^*_{\ik{N}}\big)\Big].
\end{split}
  \label{eq:decrease}
\end{equation}
The first term is the cost paid at the current step and is never positive. The whole difference
is therefore non-positive as soon as the bracket is, and the bracket involves the terminal
ingredients only.

\begin{theorem}[stability of MPC]
Let $f(0,0) = 0$, let $\ell$ and $V_f$ be positive definite, let $\Xf$ and $\kappa_f$ satisfy
\eqref{eq:terminalinv}, and let
\begin{equation}
  V_f\big(f(x, \kappa_f(x))\big) - V_f(x) \;\le\; -\,\ell\big(x, \kappa_f(x)\big)
  \qquad \text{for every } x \in \Xf .
  \label{eq:terminalcond}
\end{equation}
Then $V^*_N$ is a Lyapunov function for the closed loop, $\Delta V^*_N(x_k) \le -\ell(x_k,
\kappa_N(x_k))$, and the origin is asymptotically stable with the feasible set as region of
attraction.
\end{theorem}

Condition \eqref{eq:terminalcond} says that the terminal cost must itself be a Lyapunov
function for the system driven by the terminal law, decreasing by at least the stage cost at
every step. Stability of MPC is therefore not something one hopes for: it is obtained by
choosing $V_f$ and $\kappa_f$ properly. The practical reading is the following. If a feedback
law $u = c(x)$ is known to stabilize the origin, with a Lyapunov function $W$ proving it and a
decrease of at least the stage cost, then taking $V_f \equiv W$ and $\kappa_f \equiv c$ makes
the MPC controller stabilizing: the proof above becomes exactly the property of $W$. For the
linear-quadratic problem this is the LQR law with $V_f(x) = x\T W x$, the same choice that
\cref{sec:optimality} obtained for optimality.

\subsection{The role of each ingredient}

The theory of the last three chapters assigns a role to each ingredient of $\mathcal{P}_N(x_k)$.
\begin{itemize}
  \item The \emph{horizon} $N$ trades optimality for computation: as $N$ grows the problem
        approaches the infinite-horizon one, which is the best that can be done and which cannot
        be solved; at the same time the number of variables and constraints grows, and so does
        the solution time.
  \item The \emph{terminal cost} $V_f$ serves optimality, when it is the value function of the
        infinite-horizon problem, and stability, when it satisfies \eqref{eq:terminalcond}. When
        the value function is not available it can be approximated, or learned.
  \item The \emph{terminal set} $\Xf$ serves recursive feasibility, through
        \eqref{eq:terminalinv}. It is an artificial constraint, added for the proof and not by
        the application, and it costs performance: the larger it can be made, the better, the
        limit being $\Xf = \X$ when $\X$ itself satisfies \eqref{eq:terminalinv}.
\end{itemize}

\begin{reading}
\refkey{Cannon\_2016}, section 2.6, ``Stability and Convergence'' (pp.~29--32): the proof of this
chapter in the same form, with the optimal cost as Lyapunov function and the condition on the
terminal cost. \refkey{Rawlings\_2017}, appendix B, section B.2, ``Stability and Asymptotic
Stability'' (pp.~696--700), for the definitions, and section B.3, ``Lyapunov Stability
Theory'' (pp.~700--709), for the Lyapunov theorems in their general form; section 2.4.2,
``Stabilizing Conditions'' (pp.~114--120), for the assumptions on $\ell$, $V_f$, $\kappa_f$ and
$\Xf$ as they are used throughout that book. \refkey{Mayne\_2000} is the survey in which these
conditions were first collected in the form used today.
\end{reading}
